% KYWO SYSTEM - EXPERIMENTAL TEST SUMMARY TABLES
% Date: March 16, 2026
% Copy and paste these tables into your LaTeX document

% Test 0: Software Ping-Pong Description
% This test measures round-trip latency and inter-device synchronization using
% software timestamps. The PC broadcasts a UDP packet to both edge nodes simultaneously,
% and each node responds with an acknowledgment. Software drift is calculated as the
% timing difference between the two device responses, providing a software-based
% estimate of synchronization quality.

% Table: Test 0 - Software Ping-Pong Results (Architecture 1)
\begin{table}[htbp]
\centering
\caption{Test 0: Software Ping-Pong Round-Trip Latency and Drift}
\begin{tabular}{|l|c|c|c|c|c|c|}
\hline
\textbf{Metric} & \textbf{Mean} & \textbf{Median} & \textbf{Std Dev} & \textbf{Range} & \textbf{P95} & \textbf{P99} \\ \hline
Device A RTT (ms) & 57.03 & 41.72 & 54.97 & 4.01--257.76 & 175.23 & 226.86 \\ \hline
Device B RTT (ms) & 56.89 & 41.45 & 54.92 & 3.87--257.72 & 175.22 & 226.85 \\ \hline
Software Drift (ms) & 0.50 & 0.02 & 0.96 & 0.01--6.59 & 2.53 & 3.73 \\ \hline
\end{tabular}
\label{tab:test0_results}
\end{table}

% Test 1: Hardware Observer - Architecture 1 Description
% This test uses a hardware observer with microsecond-precision interrupts to measure
% ground truth latency and synchronization. The PC triggers the observer and broadcasts
% a UDP packet to both edge nodes. Each node responds by setting a GPIO pin HIGH, which
% the observer captures via hardware interrupts. This eliminates software timestamp
% inaccuracies and provides the true physical synchronization between devices.

% Table: Test 1 - Complete Hardware Observer Results (Architecture 1)
\begin{table}[htbp]
\centering
\caption{Test 1: Hardware Observer Latency and Synchronization (Architecture 1)}
\begin{tabular}{|l|c|c|c|c|c|c|}
\hline
\textbf{Metric} & \textbf{Mean} & \textbf{Median} & \textbf{Std Dev} & \textbf{Range} & \textbf{P95} & \textbf{P99} \\ \hline
Device A Latency (ms) & 49.42 & 25.16 & 48.68 & 2.05--210.51 & 110.18 & 188.93 \\ \hline
Device B Latency (ms) & 49.42 & 25.16 & 48.69 & 2.06--210.50 & 110.19 & 188.95 \\ \hline
Sync Drift ($\mu$s) & 7.88 & 5.00 & 10.34 & 1--99 & 21.00 & 32.50 \\ \hline
\end{tabular}
\label{tab:test1_results}
\end{table}

% Test 2: Hardware Observer - Architecture 2 Description
% This test measures synchronization drift in autonomous operation. Two edge nodes run
% independent state machines (grandmaster and follower) that execute sequences without
% centralized control. The hardware observer passively monitors GPIO timing between the
% two nodes. ESP-NOW clock synchronization occurs every 2000ms to correct for hardware
% clock drift. This test reveals the long-term stability of distributed execution.

% Table: Test 2 - Hardware Observer Results (Architecture 2)
\begin{table}[htbp]
\centering
\caption{Test 2: Synchronization Drift in Autonomous Execution (Architecture 2)}
\begin{tabular}{|l|c|c|c|c|c|c|}
\hline
\textbf{Metric} & \textbf{Mean} & \textbf{Median} & \textbf{Std Dev} & \textbf{Range} & \textbf{P95} & \textbf{P99} \\ \hline
Sync Drift ($\mu$s) & 108.47 & 31.00 & 495.96 & 1--6449 & 406.10 & 1563.89 \\ \hline
\end{tabular}
\label{tab:test2_results}
\end{table}

% Table: Comparative Analysis
\begin{table}[htbp]
\centering
\caption{Comparative Analysis: Architecture 1 vs Architecture 2 Synchronization Performance}
\begin{tabular}{|l|c|c|c|}
\hline
\textbf{Metric} & \textbf{Arch 1 (UDP)} & \textbf{Arch 2 (ESP-NOW)} & \textbf{Ratio (A2/A1)} \\ \hline
Mean Drift ($\mu$s) & 7.88 & 108.47 & 13.8$\times$ \\ \hline
Median Drift ($\mu$s) & 5.00 & 31.00 & 6.2$\times$ \\ \hline
Std Dev ($\mu$s) & 10.34 & 495.96 & 48.0$\times$ \\ \hline
Max Drift ($\mu$s) & 99 & 6449 & 65.1$\times$ \\ \hline
P95 Drift ($\mu$s) & 21.00 & 406.10 & 19.3$\times$ \\ \hline
P99 Drift ($\mu$s) & 32.50 & 1563.89 & 48.1$\times$ \\ \hline
\end{tabular}
\label{tab:comparative_analysis}
\end{table}

% Key Findings (use in document text):
% Architecture 1 achieves 13.8$\times$ better mean synchronization
% Architecture 1 has 48.0$\times$ more stable timing (lower std dev)
% Architecture 1 worst-case drift: 99 $\mu$s
% Architecture 2 worst-case drift: 6449 $\mu$s (65.1$\times$ worse)
